\subsection{Dexterous Manipulation}
Dexterous manipulation has been an active area of research for decades~\citep{DBLP:conf/icra/Fearing86, DBLP:journals/ijrr/Rus99,DBLP:journals/trob/Bicchi00, DBLP:conf/icra/OkamuraSC00, DBLP:conf/icar/MaD11}.
Many different approaches and strategies have been proposed over the years.
This includes rolling~\citep{DBLP:conf/icra/BicchiS95, DBLP:conf/icra/HanGLQT97, DBLP:conf/icra/HanT98, DBLP:journals/trob/CherifG99, DBLP:conf/icra/DoulgeriD13}, sliding~\citep{DBLP:journals/trob/CherifG99, DBLP:journals/trob/ShiWUL17}, finger gaiting~\citep{DBLP:conf/icra/HanT98}, finger tracking~\citep{DBLP:conf/icra/Rus92}, pushing~\citep{DBLP:journals/corr/DafleR17}, and re-grasping~\citep{DBLP:conf/icra/TournassoudLM87, DBLP:conf/icra/DafleRPTSEMLSF14}.
For some hand types, strategies like pivoting~\citep{DBLP:conf/iros/AiyamaII93}, tilting~\citep{DBLP:journals/trob/ErdmannM88}, tumbling~\citep{sawasaki1991tumbling}, tapping~\citep{DBLP:journals/ijrr/HuangM00}, two-point manipulation~\citep{DBLP:conf/iros/AbellE95}, and two-palm manipulation~\citep{DBLP:journals/ijrr/Erdmann98} are also options.
These approaches use planning and therefore require exact models of both the hand and object.
After computing a trajectory, the plan is typically executed open-loop, thus making these methods prone to failure if the model is not accurate.\footnote{Some methods use iterative re-planning to partially mitigate this issue.}

Other approaches take a closed-loop approach to dexterous manipulation and incorporate sensor feedback during execution, e.g. tactile sensing~\citep{DBLP:conf/icra/TaharaAY10, DBLP:conf/iros/LiBKB14, DBLP:conf/icra/LiYTB14, DBLP:journals/ijma/0001MHRB13}.
While those approaches allow for the correction of mistakes at runtime, they still require reasonable models of the robot kinematics and dynamics, which can be challenging to obtain for under-actuated hands with many degrees of freedom.

Deep reinforcement learning has also been used successfully to learn complex manipulation skills on physical robots.
Guided policy search~\citep{DBLP:conf/icml/LevineK13, DBLP:conf/icra/LevineWA15} learns simple local policies directly on the robot and distills them into a global policy represented by a neural network. Soft Actor-Critic\cite{DBLP:journals/corr/abs-1812-05905} has been recently proposed as a state-of-the-art model-free algorithm optimizing concurrently both expected reward and action entropy, that is capable of learning complex behaviors directly in the real world.

Alternative approaches include using many physical robots simultaneously, in order to be able to collect sufficient experience~\citep{DBLP:conf/icra/GuHLL17, DBLP:journals/ijrr/LevinePKIQ18, 2018arXiv180610293K} or leveraging a model-based learning algorithms, which generally possess much more favorable sample complexity characteristics~\cite{DBLP:journals/corr/abs-1808-09105}. Some researchers have successfully  utilized expert human demonstrations in guiding the training process of the agents~\citep{DBLP:journals/corr/abs-1906-11695,DBLP:journals/corr/abs-1903-01973,DBLP:journals/corr/abs-1810-01845,DBLP:journals/corr/abs-1810-06045}.

\subsection{Dexterous In-Hand Manipulation}
Since a very large body of past work on dexterous manipulation exists, we limit the more detailed discussion to setups that are most closely related to our work on dexterous in-hand manipulation.

Mordatch et al.~\citep{DBLP:conf/sca/MordatchPT12} and Bai et al.~\citep{DBLP:conf/icra/BaiL14} propose methods to generate trajectories for complex and dynamic in-hand manipulation, but their results are limited to simulation.
There has also been significant progress in learning complex in-hand dexterous manipulation~\citep{plappert2018multi, DBLP:journals/corr/abs-1804-08617}, tool use~\citep{DBLP:journals/corr/abs-1709-10087} and even solving a smaller model of a Rubik's Cube \cite{2019arXiv190711388L} using deep reinforcement learning, but those approaches were likewise only evaluated in simulation.

In contrast, multiple authors learn policies for dexterous in-hand manipulation directly on the robot.
Hoof et al.~\citep{DBLP:conf/humanoids/HoofHN015} learn in-hand manipulation for a simple 3-fingered gripper whereas Kumar et al.~\citep{DBLP:conf/icra/KumarTL16, DBLP:journals/corr/KumarGTL16} and Falco et al.~\citep{falco2018policy} learn such policies for more complex humanoid hands. In \cite{PDDM}, the authors learn a forward dynamics model and use model predictive control to manipulate two Baoding balls with a Shadow hand.
While learning directly on the robot means that modeling the system is not an issue, it also means that learning has to be performed with limited data.
This is only possible when learning simple (e.g. linear or local) policies that, in turn, do not exhibit sophisticated behaviors.

\subsection{Sim to Real Transfer}

\emph{Domain adaption} methods~\citep{DBLP:journals/corr/TzengDHFPLSD15, DBLP:journals/corr/GuptaDLAL17}, progressive nets~\citep{DBLP:conf/corl/RusuVRHPH17}, and learning inverse dynamics models~\citep{DBLP:journals/corr/ChristianoSMSBT16} were all proposed to help with sim to real transfer.
All of these methods assume access to real data.
An alternative approach is to make the policy itself more adaptive during training in simulation using \emph{domain randomization}.
Domain randomization was used to transfer object pose estimators~\citep{tobin2017domain} and vision policies for fly drones~\citep{DBLP:conf/rss/SadeghiL17}.
This idea has also been extended to dynamics randomization~\citep{peng2017sim, DBLP:journals/corr/AntonovaCSK17, DBLP:journals/corr/abs-1804-10332, DBLP:conf/rss/YuTLT17, openai2018learning} to learn a robust policy that transfers to similar environments but with different dynamics.
Domain randomization was also used to plan robust grasps~\citep{DBLP:conf/rss/MahlerLNLDLOG17, DBLP:journals/corr/abs-1709-06670, DBLP:journals/corr/abs-1710-06425} and to transfer learned locomotion~\citep{DBLP:journals/corr/abs-1804-10332} and grasping~\citep{DBLP:journals/corr/abs-1802-09564} policies for relatively simple robots.
Pinto et al.~\citep{pinto2017robust} propose to use \emph{adversarial training} to obtain more robust policies and show that it also helps with transfer to physical robots~\citep{pinto2017supervision}. Hwangbo et al.~\cite{hwangbo2019learning} used real data to learn an actuation model and combined it with domain randomization to successfully transfer locomotion policies.

A number of recent works have focused on adapting the environment distribution of domain randomization to improve sim-to-real transfer performance. For policy training, one approach viewed the problem as a bi-level optimization~\citep{Vuong2019, Ruiz2018}. Chebotar et al.~\citep{Chebotar2018} used real-world trajectories and a discrepancy metric to guide the distribution search. In~\citep{Mehta2019}, a discriminator was used to guide the distribution in simulation. For vision models, domain randomization has been modified to improve image content diversity~\citep{Kar2019, Prakash2018, Cubuk2018} and to adapt the distribution by using an adversarial network~\citep{Zakharov2019}.

\subsection{Meta-Learning via Reinforcement Learning}

Despite being a very young field, meta-learning in the context of deep reinforcement learning already has a large body of work published. Algorithms such as MAML~\cite{DBLP:journals/corr/FinnAL17} and SNAIL~\cite{DBLP:journals/corr/MishraRCA17} have been developed to improve the sample efficiency of reinforcement learning agents. A common theme in research is to try to exploit a shared structure in a distribution of environments, to quickly identify and adapt to previously unseen cases~\cite{DBLP:journals/corr/abs-1907-04964, Smundsson2018MetaRL,DBLP:journals/corr/abs-1905-06424,Lan2019MetaRL}
There are works that directly treat meta-learning as identification of a dynamics model of the environment~\cite{DBLP:journals/corr/abs-1803-11347} while others tackle the problem of task discovery for training~\cite{DBLP:journals/corr/abs-1806-04640}.
Meta-learning has also been studied in a multi-agent setting~\cite{DBLP:journals/corr/abs-1903-02710}.

The approach we've taken is directly based on $RL^2$~\cite{DBLP:journals/corr/DuanSCBSA16,DBLP:journals/corr/WangKTSLMBKB16} where a general-purpose optimization algorithm trains a model augmented with memory to perform independent learning algorithm in the inner loop. The novelty in our results comes from the combination of automated curriculum generation (ADR), a challenging underlying problem (solving Rubik's Cube) and a completely out-of-distribution test environment (sim2real). 

This approach coincides with what Jeff Clune described as AI-GA~\cite{clune2019ai}, the AI-generating algorithms, whose one of three pillars is generating effective and diverse learning environments. In similar spirit to our ADR, related works include  PowerPlay~\cite{DBLP:journals/corr/abs-1112-5309,DBLP:journals/corr/abs-1210-8385}, POET~\cite{wang2019poet} and different varieties of self-play~\cite{alphago,DBLP:journals/corr/abs-1712-01815,five}.


\subsection{Robotic Systems Solving Rubik's Cube}

While many robots capable of solving a Rubik's Cube exist today~\cite{rubikguinnes,mindcuber,rcr3d,higo2018rubik}, all of them which we are aware of have been built exclusively for this purpose and cannot generalize to other manipulation tasks.






